\section{Literature Review}
\label{literature_review}

Facility location problems (FLPs) are a broad class of optimization problems that determine the optimal positions for one or more facilities to optimize specific criteria (e.g., cost, coverage, service times), subject to constraints \cite{farahani2009facility}. Although FLPs are not uniquely classified in literature, they may be classified into four categories according to Adeleke and David work \cite{adeleke2020facility}: (1) analytical models, based on theoretical assumptions; (2) network models, which use transportation graphs to compute shortest paths; (3) continuous models, where facilities may be located anywhere on a continuous plane; and (4) discrete models, which restrict facility placement to a predefined set of candidate sites. In most practical cases, discrete models are particularly relevant due to their ability to incorporate realistic constraints, such as land availability \cite{adeleke2020facility}. 

Within the FLP group, a more specific branch is the Coverage Location Problem (CLP). Formalized by Toregas et al. \cite{toregas1971,berman2010, broin1986}, CLP ensures that demand is covered within a specific service radius. The fundamental idea behind the CLP is to determine the optimal location of a minimum number of facilities so that all points of demand are within range of at least one of them. From the CLP, the Maximum Coverage Location problem (MCLP) was introduced by Church and Revelle \cite{church1974}. In contrast to the CLP approach, which aims to cover all demand with the minimum possible number of facilities, the MCLP assumes that in many situations, it is impossible to cover all demand points. Hence, its objective is to maximize demand coverage given limited facilities. Therefore, it is often applied in contexts with limited resources, such as waste management, where only a limited number of containers are available. Both models provide a foundation for planning urban infrastructure where resources and demand must be carefully balanced.

For instance, Hartini et al. \cite{hartini2021design} designed a network of UCO collection centers in Semarang, Indonesia. Previously, the authors conducted a survey of several households to understand waste-disposal behavior. Then, they formulated an MCLP to determine the optimal number and location of collection centers for two defined distance threshold scenarios (1 and 2 km). However, their model was based on a single-objective formulation that considered only coverage, ignoring economic costs. Consequently, the study also did not address whether the distribution of collection centers was equitable. Moreover, the authors used the total number of households, the average monthly UCO generation per household, and the estimated participation rate to estimate waste disposal at each collection point. Although this approach captures the absolute volume of waste generated, it does not incorporate population density as a spatial planning criterion. As a result, the model does not explicitly account for population concentration or dispersion, which may influence accessibility and equity in service provision.
In contrast, using real-world population density data, Moreno et al. \cite{moreno2024optimizing} applied the MCLP to optimize the placement of UCO containers in Valencia, Spain. Their approach utilized a genetic algorithm (GA) to solve the problem, incorporating a feasibility function to ensure compliance with the maximum number of containers currently deployed in the city. The study demonstrated that evolutionary algorithms are well-suited for addressing the MCLP and improving the arrangement of UCO containers. While their approach effectively improved coverage by considering population density when deciding where to place UCO containers, it focused solely on a single objective, neglecting cost and spatial fairness. 

As can be seen, single-objective approaches are limited in urban planning contexts, where locating collection points involves balancing multiple factors, that are often inversely proportional \cite{farahani2010multiple}. Therefore, when planning public infrastructure, it is necessary to explicitly consider the balance by transitioning to multi-objective models. Instead of optimizing a single criterion, multi-objective approaches aim to identify a set of Pareto-optimal solutions that represent different trade-offs among the considered objectives \cite{tuvsar2014visualization}. With this paradigm shift, and considering as many objectives as deemed necessary, subsequent decision-making on the solutions obtained will take into account the inherent complexity in infrastructure planning. In that sense, Adeleke and Ali \cite{adeleke2021efficient} proposed a mathematical model to locate solid waste collection sites in Eti-Osa LGA, Lagos, Nigeria. Specifically, they formulated a set coverage problem (SCLP), another variant of the FLP, whose objective is to minimize the number of collection points while ensuring that all users are covered within a specified distance. Their model assumes that the quantities and types of waste generated per user are known. Then, a Lagrangian relaxation combined with subgradient optimization is used to solve the bi-objective problem. The computational results, based on simulated data for the Nigerian municipality, show significant reductions in the number of containers assigned, demonstrating the effectiveness of their proposal. Although their method introduced operational realism and achieved significant efficiency gains, it did not explicitly consider multi-objective trade-offs. Similarly, Rossit and Bard \cite{rossit2025solving} proposed a two-objective optimization model for the problem of waste container location under uncertain waste generation rates. Their model simultaneously minimizes costs and collection frequency, providing a flexible framework for exploring more efficient solutions. However, their approach does not include spatial equity or population coverage criteria. Likewise, their model also faced computational limitations when applied to large urban scenarios. As the number of waste generation points increased, the number of variables and constraints in the mixed-integer linear programming (MILP) formulation grew fast, making the problem computationally intensive. For example, in their largest instance with 190 generation points, the robust model involved over 1.1 million constraints and nearly 472,000 variables. Despite using advanced tools such as Gurobi, the authors reported that the solver failed to find feasible solutions within the time limit in some urban-demanding scenarios.  

%\subsection{Multi-Objective Evolutionary Algorithms for FLP}
Given these computational challenges, attention has shifted toward heuristic and metaheuristic methods better suited to exploring large and complex solution spaces. In particular, this computational complexity has led to the adoption of Multi-objective Evolutionary Algorithms (MOEAs), with the non-dominated sorting genetic algorithm II (NSGA-II) being especially popular due to its efficiency. For example, Harris et al. \cite{harris2009multi} extended the classic Uncapacitated Facility Location Problem (UFLP), traditionally formulated as a single-objective problem, to a multi-objective model that integrates economic and environmental criteria within a green logistics framework. Their study represents one of the first efforts to go beyond cost-centered formulations, demonstrating that NGSA-II can generate Pareto fronts that capture trade-offs between operational efficiency and environmental sustainability. By modeling the environmental impacts of transportation and depot operations separately, they demonstrated how MOEAs can support strategic facility planning that addresses trade-offs between conflicting sustainability goals in infrastructure design.

However, the performance of NSGA-II deteriorates as the number of objectives increases \cite{ishibuchi2016performance}. This situation led to the introduction of NSGA-III, which, using a reference-point-based diversity preservation strategy, has demonstrated superior performance in applications with three or more objectives \cite{vargas2025resilient}. In this regard, Jain and Deb \cite{deb2013evolutionary, jain2013evolutionary} demonstrated how the new version of the algorithm shows better diversity and convergence than the previous version in high-dimensional spaces. With three objectives, as in this work, NSGA-II remains a standard and competitive choice \cite{ishibuchi2016performance}, and its crowding distance is simpler to adapt to fronts whose first objective takes few discrete values. 

\paragraph{Spatially aware genetic operators.}
The genetic operators used in a GA play a fundamental role in the effectiveness of evolutionary algorithms, especially in spatial optimization. Specifically, traditional crossover operators (e.g., uniform crossover) treat genes independently, without accounting for possible geographical relationships between them. Therefore, spatially coherent solutions may be destroyed during recombination, leading to lower-quality solutions. Based on this idea, some authors have developed more effective operators by considering various geographical relationships between genes. For example, Li and Parrott \cite{li2016improved} proposed an implementation of a GA to solve a multi-site land use allocation (MLUA) problem, an optimization problem that attempts to determine how different parcels of land across multiple sites should be allocated to various land uses (e.g., residential, industrial, agricultural, etc.) in order to achieve specific objectives (e.g., maximizing economic return, minimizing environmental impact, etc.). In this context, their GA incorporated two spatially aware crossovers that improved the algorithm's performance. Their first one considered local spatial relationships by replacing a cell in one parent with the corresponding cell from another if the surrounding neighborhood supported the same land use type, thereby reducing fragmentation. The second one introduced global diversity by swapping quadrants of the spatial grid between parents, promoting the exploration of the solution space.
Additionally, two mutation operators were introduced. The first mutates entire clusters of adjacent cells rather than individual cells, maintaining spatial contiguity. The second one promotes compactness by probabilistically converting a cell and its neighborhood based on the number of adjacent cells with the same land-use type.

Liu and Cho \cite{liu2020spatially} developed a Spatially Explicit Evolutionary Computation (SEEC) framework for spatial partitioning problems. Their approach replaces traditional recombination with graph-based operators that explicitly preserve spatial contiguity and topology. They introduced two novel operators. The first is a path-relinking-based crossover (PRCRX) that constructs a path between two parent solutions by incrementally transforming one into the other through spatially valid moves. The second is an ejection-chain-based mutation (ECMUT) that performs chained local mutations, maintaining contiguity. Both are designed to avoid infeasible solutions, significantly improving solution quality and outperforming conventional heuristics. Similarly, a recent work by Jordan et al. \cite{jordan2025crossover} introduced a graph-based crossover operator for problems where the objective can be defined over a graph of neighborhoods. Their crossover operator exchanges substructures between parents by selecting and swapping neighborhoods, preserving local interdependencies and structural coherence by inheriting entire subgraphs defined by a center node and its neighbors up to a configurable depth. The operator is general-purpose, independent of the specific graph topology, and has been shown to outperform or match specialized crossover operators across most of the tested problems, particularly in scenarios with strong variable interdependencies. In this paper we test this principle on a location problem with two spatially structured crossovers, the graph crossover and a regional crossover that exchanges the candidate locations nearest to a random centre, and find that, in our problem, recombination brings no benefit over a well-designed mutation.

Overall, while significant progress has been made in optimizing the placement of waste collection infrastructure, most approaches address only a single objective or do not explicitly account for spatial equity. Moreover, the reviewed literature emphasizes the importance of spatially informed genetic operators for addressing complex optimization problems in which variables exhibit spatial interdependence. 
In the following section, we formally define the problem of UCO collection point placement, extending classical FLP and MCLP formulations to incorporate cost, population coverage, and spatial equity considerations.